% ============================================================
% preamble.tex — paket, gaya, dan makro dokumentasi TUSKO
% ============================================================
\usepackage[T1]{fontenc}
\usepackage[utf8]{inputenc}
\usepackage{lmodern}
\usepackage[a4paper,margin=2.5cm,includefoot]{geometry}
\usepackage{graphicx}
\usepackage{xcolor}
\usepackage{float}
\usepackage{fancyhdr}
\usepackage{titlesec}
\usepackage{enumitem}
\usepackage{booktabs}
\usepackage{longtable}
\usepackage{array}
\usepackage{tcolorbox}
\usepackage{caption}
\usepackage{hyperref}
\usepackage{microtype}
\usepackage{ifthen}

\definecolor{tuskoamber}{HTML}{D97706}
\definecolor{tuskodark}{HTML}{171717}
\definecolor{tuskogray}{HTML}{737373}
\definecolor{tuskolight}{HTML}{F5F5F5}

\hypersetup{
  colorlinks=true,
  linkcolor=tuskodark,
  urlcolor=tuskoamber,
  citecolor=tuskodark,
  pdftitle={Panduan Pengguna TUSKO},
  pdfauthor={Tim TUSKO}
}

% --- Judul otomatis Bahasa Indonesia
\renewcommand{\contentsname}{Daftar Isi}
\renewcommand{\figurename}{Gambar}
\renewcommand{\tablename}{Tabel}
\renewcommand{\refname}{Referensi}

% --- Header / footer
\pagestyle{fancy}
\fancyhf{}
\fancyhead[L]{\footnotesize\color{tuskogray}TUSKO --- Panduan Pengguna}
\fancyhead[R]{\footnotesize\color{tuskogray}\leftmark}
\fancyfoot[C]{\footnotesize\thepage}
\renewcommand{\headrulewidth}{0.4pt}

% --- Gaya judul
\titleformat{\section}{\Large\bfseries\color{tuskodark}}{}{0pt}{}
\titleformat{\subsection}{\large\bfseries\color{tuskodark}}{}{0pt}{}
\titlespacing*{\section}{0pt}{1.4em}{0.6em}

% --- \flowmeta{KODE}{Judul}{Aktor}
\newcommand{\flowmeta}[3]{%
  \section{#2}\label{flow:#1}%
  \noindent\textbf{Kode Alur:} \texttt{#1}\quad\textbar\quad\textbf{Aktor:} #3\\[0.2em]
  \noindent{\color{tuskoamber}\rule{\linewidth}{1.2pt}}\vspace{0.4em}%
}

% --- Lingkungan seksi flow
\newenvironment{flowsection}[1]{\subsection{#1}}{\par\vspace{0.3em}}

% --- \shot{path}{caption}{label} : gambar bila ada, else placeholder
\newcommand{\shot}[3]{%
  \begin{figure}[H]\centering%
    \IfFileExists{#1}{%
      \includegraphics[width=0.95\linewidth]{#1}%
    }{%
      \fbox{\begin{minipage}[c][0.30\textheight][c]{0.90\linewidth}\centering%
        \color{tuskogray}\textbf{[GAMBAR BELUM TERSEDIA]}\\[0.6em]%
        {\ttfamily\small #1}\\[0.6em]%
        {\itshape\small #2}\end{minipage}}%
    }%
    \caption{#2}\label{#3}%
  \end{figure}%
}

% --- \docver{versi}{tanggal}
\newcommand{\docver}[2]{\par\vspace{0.5em}\noindent{\footnotesize\itshape Dokumen #1 --- #2}\par\vspace{0.5em}}
